\documentclass[11pt,a4paper]{article}
\usepackage[margin=2.9cm]{geometry}
\usepackage{amsmath,amssymb,amsthm,mathrsfs}
\usepackage[T1]{fontenc}
\usepackage{microtype}
\usepackage{booktabs,longtable,array}
\usepackage{bm}
\usepackage[colorlinks=true,linkcolor=blue,citecolor=blue]{hyperref}

% ------------------------------------------------------------------
% Phase V1 of the gate-(II.11) execution (scoping: gate_II11_scoping.tex).
% Contents: Statements A2(a) (lattice trace representation with a marked
% slice, incl. the exact mean-shift/jump representation) and A3 (free
% thermal Weyl functionals, lattice and continuum, and the thermal
% fixed-coarse rate), all with complete proofs.  A2(b)-interacting is
% stated with its precise input clause and deferred to V3 (amendment
% recorded in Sec. 7).  Acceptance: the closed forms proved here are the
% formulas verified numerically in v0_checkpoint.py (V0 tables).
% ------------------------------------------------------------------

\newcolumntype{P}[1]{>{\raggedright\arraybackslash}p{#1}}

\theoremstyle{plain}
\newtheorem{theorem}{Theorem}[section]
\newtheorem{proposition}[theorem]{Proposition}
\newtheorem{lemma}[theorem]{Lemma}
\newtheorem{corollary}[theorem]{Corollary}
\theoremstyle{definition}
\newtheorem{definition}[theorem]{Definition}
\newtheorem{remark}[theorem]{Remark}

\newcommand{\R}{\mathbb{R}}
\newcommand{\Z}{\mathbb{Z}}
\newcommand{\C}{\mathbb{C}}
\newcommand{\T}{\mathbb{T}}
\newcommand{\E}{\mathbb{E}}
\newcommand{\W}{\mathfrak{W}}
\newcommand{\HH}{\mathcal{H}}
\newcommand{\Om}{\Omega}
\newcommand{\eps}{\varepsilon}
\newcommand{\muv}{\mu}
\newcommand{\ip}[2]{\langle #1,#2\rangle}
\DeclareMathOperator{\supp}{supp}
\DeclareMathOperator{\Trace}{Tr}
\DeclareMathOperator{\ch}{cosh}
\DeclareMathOperator{\sh}{sinh}
\DeclareMathOperator{\cth}{coth}

\title{Gate (II.11), phase V1:\\[4pt]
\large Trace representations with a marked slice, the free thermal Weyl
functionals,\\ and the thermal fixed-coarse rate}
\author{}
\date{12 August 2026}

\begin{document}
\maketitle

\begin{abstract}
This document executes phase V1 of the thermal two-torus route
(\texttt{gate\_II11\_}\allowbreak\texttt{scoping.tex}): Statements A2(a) and A3 are proved in full,
and A2(b) is put into its final conditional form. Contents: the Schr\"odinger
realization of the lattice Weyl operators with the exact CCR phase; trace-class
bounds and eigenfunction estimates giving the kernel identity for
$\Trace(e^{-tK_{N,L}}W_N(\xi))$ with complete Fubini justification; the exact
\emph{mean-shift representation} of the marked-slice trace --- the momentum
component of $\xi$ enters through an explicit harmonic jump interpolant $H_p$
which satisfies the discrete Mehler-chain harmonicity \emph{exactly} at every
time discretization, so the Gaussian change of variables is exact, its constant
is the classical action of the interpolant --- precisely the momentum half of
the thermal covariance form --- and the shift phase cancels the CCR phase
identically; the closed free thermal Weyl functionals on both sides
(lattice: finite mode products; continuum: countable products with a complete
convergence argument), by an elementary Laguerre-series proof of the
single-mode thermal displacement formula; and the thermal fixed-coarse rate
theorem, uniform on $t\in[t_0,\infty]$, obtained from the proved vacuum-rate
lemmas with two new derivative bounds and explicit constants. Every import is
isolated with its hypotheses; the V0 numeric tables verify each closed form
proved here.
\end{abstract}

\tableofcontents

\section{Conventions, imports, and the proof contract}\label{sec:conv}

We work in the setting of the Part II model sheet
(\texttt{lamport\_part\_ii\_}\allowbreak\texttt{model\_sheet.tex}), whose
conventions are used without change: the spatial torus $\T_L=\R/(2L\Z)$; the
lattice $\Lambda_N$ of $2r_N$ sites with spacing $\eps_N=L/r_N$ and momenta
$\Gamma_N=\{\pi j/L:-r_N\le j\le r_N-1\}$; the dispersion
$\gamma_N(k)=(m^2+4\eps_N^{-2}\sin^2(\eps_Nk/2))^{1/2}$; the canonical pairs
$[\Phi_N(x),\Pi_N(y)]=i\eps_N^{-1}\delta_{xy}$; the smeared fields
$\Phi_N(q)=\eps_N\sum_xq(x)\Phi_N(x)$, $\Pi_N(p)=\eps_N\sum_xp(x)\Pi_N(x)$; the
Weyl operators $W_N(\eps_Nq+ip)=e^{i(\Phi_N(q)+\Pi_N(p))}$; the hat
$\widehat f(k)=\eps_N^{1/2}\sum_xf(x)e^{-ikx}$; the free Hamiltonian
$H^{(0)}_{N,L}$ of model-sheet (3.1); the Wick covariance $c_{N,L}$ of (3.4);
and the interacting operator
\[
 K_{N,L}=H^{(0)}_{N,L}+V_N,\qquad
 V_N:=\lambda\,\eps_N\!\!\sum_{x\in\Lambda_N}\!
 {:}\Phi_N(x)^4{:}_{c_{N,L}},\qquad\lambda\ge0,
\]
with its coercive closed-form realization, compact resolvent, and simple ground
state (model-sheet Proposition, \S3). Everything below except \S\ref{sec:cont}
is finite-dimensional in the number of sites (though infinite-dimensional as
quantum mechanics), and no constructive-QFT input is used.

\paragraph{Imports.} (I1) The \emph{harmonic-reference} Feynman--Kac formula
for $e^{-tK_{N,L}}$: jointly continuous, strictly positive kernel
\begin{equation}
 e^{-tK_{N,L}}(\bm x,\bm y)
 =k^{(0)}_t(\bm x,\bm y)\;
 \E^{\,t}_{\bm x,\bm y}\Bigl[\exp\Bigl(-\int_0^tV_N(B_s)\,ds\Bigr)\Bigr],
 \label{eq:FK}
\end{equation}
where $k^{(0)}_t$ is the kernel of $e^{-tH^{(0)}_{N,L}}$ (the Mehler product
of Lemma~\ref{lem:mehler}), $\E^t_{\bm x,\bm y}$ the corresponding
\emph{oscillator}-bridge expectation, and $V_N$ the \emph{interaction} alone
--- not the full multiplication potential: the harmonic part is carried by the
reference kernel. Proof: V3 (\texttt{v3\_continuum\_package.tex}), Theorem
3.1(ii), applied with the lattice dictionary of its Remark 3.2 (dispersions
$\gamma_N(k)$, $k\in\Gamma_N$; interaction $V_N$; class membership by the
model-sheet \S3 proposition) --- the Trotter/monotone-convergence argument
there is run against the Mehler reference. \emph{Two bookkeeping remarks.}
(a) V3, Theorem 3.1(ii) states the identity in the \emph{Gaussian}
($\muv$-) representation, with the normalized Mehler operator
$M_t$ ($M_t\mathbf 1=\mathbf 1$); \eqref{eq:FK} is its Lebesgue form, obtained
by the ground-state (Doob) conjugation
\[
 k^{(0)}_t(\bm x,\bm y)
 =e^{-tE^{(0)}_N}\,\Omega_0(\bm x)\,\Omega_0(\bm y)\,M_t(\bm x,\bm y),
 \qquad E^{(0)}_N=\tfrac12\sum_{k\in\Gamma_N}\gamma_N(k),
\]
$\Omega_0$ the free ground state; the two bridge measures induced on paths
coincide, because the interior factors $\Omega_0^2$ are precisely the
reference density, and the boundary factors cancel in every ratio used below.
(Lemma~\ref{lem:mehler} is stated in the rescaled coordinates
$\bm X=\eps_N^{1/2}\bm\varphi$; as a $d\bm\varphi$-kernel it carries the
constant Jacobian $\eps_N^{r_N}$, which likewise cancels in ratios.)
(b) \emph{Acyclicity.} V3's Theorem 3.1(ii) is proved from the Mehler kernel,
the elementary Trotter product formula for a bounded perturbation, and
monotone convergence; its only contact with this document is
Lemma~\ref{lem:mehler}, whose own proof (heat equation plus an $L^2$
uniqueness argument) uses nothing from (I1), and which V3 re-verifies inline
from the Hermite eigenvalues. The dependency order is therefore
Lemma~\ref{lem:mehler} $\to$ V3~3.1(ii) $\to$
\S\ref{sec:kernel} (Lemma~\ref{lem:efbound}, Theorem~\ref{thm:kernel}),
Lemmas~\ref{lem:interp}--\ref{lem:shift}, Theorem~\ref{thm:meanshift}
$\to$ V3~3.1(v)--(vi); no cycle arises. The model sheet's own formula (3.16) --- flat Gaussian
kernel, Brownian bridge, \emph{full} potential --- is a different identity and
is \emph{not} the one used below. (I2) From the free-rate document
(\texttt{lamport\_part\_ii\_free\_}\allowbreak\texttt{fixed\_coarse\_rate.tex}): the uniform $D4$
envelope (its Lemma 2.1), the filter-tail bound (Lemma 2.2), the dispersion
estimates (Lemma 3.1), its pulled-back covariances (Lemmas 1.1 \emph{and} 1.2),
and --- in \S\ref{sec:rate} --- the \emph{proof architecture} of its Theorem 4.1:
the constants $r_*,A_q,A_p,b,\kappa,c_0$ of (4.1)--(4.6) together with $C_D$ of (2.3)
and $J_r$ of (4.10), the split (4.10)--(4.11), the counting bounds (4.16) and the
integral-comparison tail (4.20)--(4.21), which \S\ref{sec:rate} re-runs verbatim with thermal weights.
(I0) The model sheet, including its \S3 Proposition (class membership of the
finite-torus systems), and the scoping note's Lemma 2.1 where cited.
(I3) Elementary classical facts proved inline where used (Mehler kernel:
verified, not imported; Laguerre summation: proved). No other input enters.

\paragraph{Contract.} Every limit names its topology; every interchange of sum
and integral is dominated explicitly; statements consumed from V0 are numerical
verifications only and never premises.

\section{The Schr\"odinger realization and the CCR phase}\label{sec:schrod}

On $\HH_N=L^2(\R^{\Lambda_N},d\bm\varphi)$ the model sheet realizes
$\Phi_N(x)$ as multiplication by $\varphi_x$ and $\Pi_N(x)=-i\eps_N^{-1}
\partial_{\varphi_x}$. Then
$\Pi_N(p)=-i\sum_xp(x)\partial_{\varphi_x}$.

\begin{proposition}[translations and the exact phase]\label{prop:phase}
\emph{(i)} $\Phi_N(q)+\Pi_N(p)$ is essentially self-adjoint on
$C_c^\infty(\R^{\Lambda_N})$, and
\begin{equation}
 \bigl(W_N(\eps_Nq+ip)\,\psi\bigr)(\bm\varphi)
 =e^{\,i\eps_Nq\cdot\bm\varphi+\frac i2\eps_N\,q\cdot p}\;
 \psi(\bm\varphi+p),
 \label{eq:Waction}
\end{equation}
where $q\cdot\bm\varphi:=\sum_xq(x)\varphi_x$ etc.
\emph{(ii)} Consequently
\begin{equation}
 W_N(\eps_Nq+ip)=e^{\frac i2\eps_N q\cdot p}\;
 e^{i\Phi_N(q)}\,e^{i\Pi_N(p)},
 \qquad
 \bigl(e^{i\Pi_N(p)}\psi\bigr)(\bm\varphi)=\psi(\bm\varphi+p).
 \label{eq:ccrphase}
\end{equation}
\end{proposition}

\begin{proof}
Define, for $s\in\R$,
\[
 (V(s)\psi)(\bm\varphi)
 :=e^{\,is\,\eps_Nq\cdot\bm\varphi+\frac i2s^2\eps_N q\cdot p}\,
 \psi(\bm\varphi+sp).
\]
Each $V(s)$ is unitary, $V(s+s')=V(s)V(s')$ by direct substitution
(the cross terms $ss'$ combine as
$s'\eps_Nq\cdot(\bm\varphi)+s\eps_Nq\cdot(\bm\varphi+s'p)
+\tfrac{\eps_N}2(s^2+s'^2)q\cdot p
=(s+s')\eps_Nq\cdot\bm\varphi+\tfrac{\eps_N}2(s+s')^2q\cdot p$),
and $s\mapsto V(s)\psi$ is continuous for $\psi\in C_c^\infty$, hence, by the
group property and density, strongly continuous everywhere. For
$\psi\in C_c^\infty$,
\[
 \frac d{ds}\Big|_{s=0}V(s)\psi
 =\bigl(i\eps_Nq\cdot\bm\varphi+p\cdot\nabla\bigr)\psi
 =i\bigl(\Phi_N(q)+\Pi_N(p)\bigr)\psi .
\]
By Stone's theorem $V$ has a self-adjoint generator $G$ extending
$(\Phi_N(q)+\Pi_N(p))|_{C_c^\infty}$; since $V(s)$ maps $C_c^\infty$ into
itself, $C_c^\infty$ is a core for $G$ (Nelson's invariant-domain criterion),
i.e.\ $\Phi_N(q)+\Pi_N(p)$ is essentially self-adjoint there and
$e^{i(\Phi_N(q)+\Pi_N(p))}=V(1)$, which is \eqref{eq:Waction}. Setting $q=0$
gives the translation formula; comparing \eqref{eq:Waction} with the
composition $e^{i\Phi_N(q)}e^{i\Pi_N(p)}$ (multiplication, then translation)
gives \eqref{eq:ccrphase}.
\end{proof}

\begin{remark}
The phase $\theta(q,p)=\tfrac12\eps_Nq\cdot p$ and the Weyl relation with the
model-sheet symplectic form were verified numerically on truncation-safe
corners in V0, with residuals $3.9\times10^{-15}$ (CCR phase) and
$2.7\times10^{-15}$ (Weyl relation).
\end{remark}

\section{Trace class, eigenfunction bounds, and the kernel identity}
\label{sec:kernel}

Throughout this section $t>0$ and $K:=K_{N,L}$ with eigenvalues
$E_0<E_1\le E_2\le\cdots\to\infty$ (compact resolvent; simplicity of $E_0$
per the model sheet) and real orthonormal eigenfunctions $\psi_j$.

\begin{proposition}[trace class]\label{prop:traceclass}
\emph{(i)} The pointwise bound
${:}u^4{:}_c=u^4-6cu^2+3c^2=(u^2-3c)^2-6c^2\ge-6c^2$ gives
$V_N\ge-12\lambda L\,c_{N,L}^2=:-c_V$.
\emph{(ii)} $E_j(K)\ge E_j(H^{(0)})-c_V$ for all $j$, and hence
\[
 Z_N(t):=\Trace e^{-tK}\le e^{tc_V}\,\Trace e^{-tH^{(0)}}
 =e^{tc_V}\prod_{k\in\Gamma_N}\bigl(2\sh(t\gamma_N(k)/2)\bigr)^{-1}<\infty .
\]
\end{proposition}

\begin{proof}
(i) is the displayed algebra summed over the $2r_N$ sites with weight
$\lambda\eps_N$ ($\eps_N\cdot2r_N=2L$). (ii): $K\ge H^{(0)}-c_V$ as quadratic
forms on the common form core, so the min--max values obey the same
inequality; both operators have discrete spectrum, and the free trace is the
standard oscillator product (computed independently in
Theorem~\ref{thm:lattice-free} below with its zero-point factors).
\end{proof}

\begin{lemma}[eigenfunction bounds]\label{lem:efbound}
For every $s>0$ and $\bm\varphi$,
\[
 |\psi_j(\bm\varphi)|\le e^{sE_j}\,
 \bigl(e^{-2sK}(\bm\varphi,\bm\varphi)\bigr)^{1/2},
 \qquad
 \int e^{-2sK}(\bm\varphi,\bm\varphi)\,d\bm\varphi=\Trace e^{-2sK}<\infty .
\]
\end{lemma}

\begin{proof}
From $\psi_j=e^{sE_j}e^{-sK}\psi_j$ and Cauchy--Schwarz,
\[
 |\psi_j(\bm\varphi)|\le e^{sE_j}\|e^{-sK}(\bm\varphi,\cdot)\|_2
 =e^{sE_j}\bigl(e^{-2sK}(\bm\varphi,\bm\varphi)\bigr)^{1/2},
\]
the last step by the semigroup property and symmetry of the (real) kernel. For the integral:
$e^{-2sK}=(e^{-sK})^2$ with $e^{-sK}$ Hilbert--Schmidt
(Proposition~\ref{prop:traceclass}) and continuous kernel \eqref{eq:FK}, so
\[
 e^{-2sK}(\bm\varphi,\bm\varphi)=\int|e^{-sK}(\bm\varphi,\bm y)|^2d\bm y,
 \qquad
 \int e^{-2sK}(\bm\varphi,\bm\varphi)\,d\bm\varphi
 \overset{\rm Tonelli}{=}\|e^{-sK}\|_{\rm HS}^2
 =\Trace e^{-2sK} .
\]
\end{proof}

\begin{theorem}[kernel identity for the marked-slice trace]\label{thm:kernel}
For all real $q,p$ and $t>0$, with $\theta=\tfrac12\eps_Nq\cdot p$,
\begin{equation}
 \Trace\bigl(e^{-tK}\,W_N(\eps_Nq+ip)\bigr)
 =e^{i\theta}\int_{\R^{\Lambda_N}}
 e^{-tK}(\bm\varphi+p,\bm\varphi)\;
 e^{\,i\eps_N q\cdot\bm\varphi}\;d\bm\varphi ,
 \label{eq:kernelid}
\end{equation}
both sides absolutely convergent.
\end{theorem}

\begin{proof}
By the spectral decomposition and boundedness of $W:=W_N(\eps_Nq+ip)$,
$\Trace(e^{-tK}W)=\sum_je^{-tE_j}\ip{\psi_j}{W\psi_j}$, absolutely convergent
since $|\ip{\psi_j}{W\psi_j}|\le1$. By \eqref{eq:Waction},
\[
 \ip{\psi_j}{W\psi_j}
 =e^{i\theta}\int\psi_j(\bm\varphi)\,
 e^{i\eps_Nq\cdot\bm\varphi}\,\psi_j(\bm\varphi+p)\,d\bm\varphi .
\]
To interchange $\sum_j$ and $\int$, dominate using
Lemma~\ref{lem:efbound} with $s=t/4$:
\[
 \sum_je^{-tE_j}|\psi_j(\bm\varphi)||\psi_j(\bm\varphi+p)|
 \le\Bigl(\sum_je^{-tE_j/2}\Bigr)
 D(\bm\varphi)^{1/2}D(\bm\varphi+p)^{1/2},
 \qquad D:=e^{-tK/2}(\cdot,\cdot)\big|_{\rm diag},
\]
and $\int D(\bm\varphi)^{1/2}D(\bm\varphi+p)^{1/2}d\bm\varphi
\le(\int D)^{1/2}(\int D(\cdot+p))^{1/2}=\Trace e^{-tK/2}<\infty$ by
Cauchy--Schwarz and Lemma~\ref{lem:efbound}. Fubini then gives
\eqref{eq:kernelid} with
$e^{-tK}(\bm x,\bm y)=\sum_je^{-tE_j}\psi_j(\bm x)\psi_j(\bm y)$, which agrees
with the continuous kernel \eqref{eq:FK} almost everywhere, hence (both sides
being continuous in $(\bm x,\bm y)$ --- the series converges locally uniformly
by the same domination) everywhere.
\end{proof}

\section{The free thermal Weyl functionals (A3(a), lattice)}\label{sec:free}

\subsection{Mode diagonalization}

Introduce the rescaled pairs $X_x:=\eps_N^{1/2}\Phi_N(x)$,
$P_x:=\eps_N^{1/2}\Pi_N(x)$, so $[X_x,P_y]=i\delta_{xy}$, and the unitary
discrete Fourier transform
$\widetilde f(k):=(2r_N)^{-1/2}\sum_xf(x)e^{-ikx}$. With
$\widetilde X(k),\widetilde P(k)$ defined likewise (so
$\widetilde X(k)^*=\widetilde X(-k)$),
the model-sheet free Hamiltonian becomes
\[
 H^{(0)}=\tfrac12\sum_{k\in\Gamma_N}
 \bigl(\widetilde P(k)^*\widetilde P(k)
 +\gamma_N(k)^2\,\widetilde X(k)^*\widetilde X(k)\bigr)
 =\sum_{k\in\Gamma_N}\gamma_N(k)\bigl(a_k^*a_k+\tfrac12\bigr),
\]
where
$a_k:=(2\gamma_N(k))^{-1/2}(\gamma_N(k)\widetilde X(k)+i\widetilde P(k))$
satisfy $[a_k,a_{k'}^*]=\delta_{kk'}$, $[a_k,a_{k'}]=0$ (direct computation
from $[\widetilde X(k),\widetilde P(k')^*]=i\delta_{kk'}$), and
\begin{equation}
 \widetilde X(k)=(2\gamma_N(k))^{-1/2}(a_k+a_{-k}^*),\qquad
 \widetilde P(k)=-i\,(\gamma_N(k)/2)^{1/2}(a_k-a_{-k}^*).
 \label{eq:modes}
\end{equation}

\begin{lemma}[displacement amplitudes]\label{lem:amp}
For real $q,p$,
$\Phi_N(q)+\Pi_N(p)=\sum_{k\in\Gamma_N}(\bar z_ka_k+z_ka_k^*)$ with
\[
\begin{gathered}
 z_k=\eps_N^{1/2}\Bigl[(2\gamma_N(k))^{-1/2}\,\widetilde q(k)
 +i(\gamma_N(k)/2)^{1/2}\,\widetilde p(k)\Bigr],\\
 \sum_{k\in\Gamma_N}|z_k|^2
 =\frac1{2\cdot2r_N}\sum_{k\in\Gamma_N}
 \Bigl(\frac{|\widehat q(k)|^2}{\gamma_N(k)}
 +\gamma_N(k)|\widehat p(k)|^2\Bigr).
\end{gathered}
\]
\end{lemma}

\begin{proof}
$\Phi_N(q)+\Pi_N(p)=\eps_N^{1/2}\sum_x(q_xX_x+p_xP_x)
=\eps_N^{1/2}\sum_k(\overline{\widetilde q(k)}\widetilde X(k)
+\overline{\widetilde p(k)}\widetilde P(k))$ by Parseval (reality of $q,p$
gives $\widetilde q(-k)=\overline{\widetilde q(k)}$). Substituting
\eqref{eq:modes} and collecting the coefficient of $a_k$ (using the $k\to-k$
symmetry of the sums) gives the displayed $z_k$. Then
\[
 |z_k|^2=\eps_N\Bigl[\frac{|\widetilde q(k)|^2}{2\gamma_N(k)}
 +\frac{\gamma_N(k)}2|\widetilde p(k)|^2
 +\operatorname{Im}\bigl(\widetilde q(k)\overline{\widetilde p(k)}\bigr)\Bigr],
\]
and the imaginary part is odd under $k\to-k$
($\widetilde q(-k)\overline{\widetilde p(-k)}
=\overline{\widetilde q(k)}\widetilde p(k)$), so it cancels in the sum;
finally $\eps_N|\widetilde f(k)|^2=|\widehat f(k)|^2/(2r_N)$.
\end{proof}

\subsection{The single-mode thermal displacement formula}

\begin{lemma}[Laguerre summation]\label{lem:laguerre}
For one mode ($[a,a^*]=1$, number basis $|n\rangle$), $z\in\C$, and
$0<w<1$:
\emph{(i)} $\ip{n}{e^{za^*}e^{-\bar za}\,n}=L_n(|z|^2)
:=\sum_{j=0}^n\binom nj\frac{(-|z|^2)^j}{j!}$;
\emph{(ii)} $\sum_{n\ge0}w^nL_n(x)=(1-w)^{-1}e^{-xw/(1-w)}$, absolutely
convergent;
\emph{(iii)} for $H=\gamma(a^*a+\tfrac12)$ and the displacement
$D(z)=e^{za^*-\bar za}$,
\[
 \frac{\Trace\bigl(e^{-tH}D(z)\bigr)}{\Trace e^{-tH}}
 =\exp\Bigl[-\tfrac12|z|^2\cth\bigl(\tfrac{t\gamma}2\bigr)\Bigr].
\]
\end{lemma}

\begin{proof}
(i) $e^{-\bar za}|n\rangle=\sum_{j\le n}\frac{(-\bar z)^j}{j!}
\sqrt{\tfrac{n!}{(n-j)!}}\,|n-j\rangle$ and likewise for the bra with $z$;
orthonormality gives
$\sum_j\frac{(-|z|^2)^j}{(j!)^2}\frac{n!}{(n-j)!}=L_n(|z|^2)$.
(ii) Insert (i) and interchange the (absolutely convergent: majorant
$\sum_j\frac{x^j}{j!}\sum_n\binom njw^n<\infty$) double sum:
\[
 \sum_{n}w^n\sum_j\binom nj\frac{(-x)^j}{j!}
 =\sum_j\frac{(-x)^j}{j!}\frac{w^j}{(1-w)^{j+1}}
 =\frac1{1-w}\,e^{-xw/(1-w)},
\]
using $\sum_{n\ge j}\binom njw^n=w^j(1-w)^{-j-1}$.
(iii) $D(z)=e^{-|z|^2/2}e^{za^*}e^{-\bar za}$ (BCH with central
$[za^*,-\bar za]=|z|^2$), so with $w=e^{-t\gamma}$ (zero-point factors cancel
in the ratio),
\[
 \frac{\Trace(e^{-tH}D(z))}{\Trace e^{-tH}}
 =e^{-|z|^2/2}\,(1-w)\sum_nw^nL_n(|z|^2)
 =e^{-|z|^2(\frac12+\frac w{1-w})}
 =e^{-\frac{|z|^2}2\frac{1+w}{1-w}},
\]
and $(1+w)/(1-w)=\cth(t\gamma/2)$.
\end{proof}

\subsection{The lattice closed form}

\begin{theorem}[A3(a), lattice]\label{thm:lattice-free}
For every $t\in(0,\infty]$ and $\xi=\eps_Nq+ip$,
\begin{equation}
 F^{(0)}_N(t;\xi):=
 \frac{\Trace\bigl(e^{-tH^{(0)}_{N,L}}W_N(\xi)\bigr)}
      {\Trace\bigl(e^{-tH^{(0)}_{N,L}}\bigr)}
 =\exp\Bigl[-\tfrac14\,Q^{(t)}_N(\xi)\Bigr],
 \label{eq:latticefree}
\end{equation}
\begin{equation}
 Q^{(t)}_N(\xi)=\frac1{2r_N}\sum_{k\in\Gamma_N}
 \cth\Bigl(\frac{t\gamma_N(k)}2\Bigr)
 \Bigl(\gamma_N(k)^{-1}|\widehat q(k)|^2
 +\gamma_N(k)|\widehat p(k)|^2\Bigr),
 \label{eq:Qt}
\end{equation}
with $\cth(\infty):=1$; moreover
$\Trace e^{-tH^{(0)}}=\prod_k(2\sh(t\gamma_N(k)/2))^{-1}$, and $t=\infty$
reproduces the model-sheet vacuum functional (3.3).
\end{theorem}

\begin{proof}
$\HH_N$ factorizes unitarily over the modes,
$H^{(0)}=\sum_k\gamma_k(a_k^*a_k+\tfrac12)$ with commuting summands, and by
Lemma~\ref{lem:amp} and the multimode BCH (all commutators central),
$W_N(\xi)=\prod_kD_k(iz_k)$ with $|iz_k|=|z_k|$. The trace factorizes over the
finite tensor decomposition; Lemma~\ref{lem:laguerre}(iii) applies per mode:
\[
 F^{(0)}_N(t;\xi)=\prod_{k\in\Gamma_N}
 e^{-\frac12|z_k|^2\cth(t\gamma_k/2)} ,
\]
and inserting Lemma~\ref{lem:amp}'s sum (with the $\cth$ weight carried along;
the cross terms cancel in the same $k\to-k$ pairing because
$\gamma_N(-k)=\gamma_N(k)$ and hence the weights are even) gives
\eqref{eq:latticefree}--\eqref{eq:Qt}. The partition function is the product
of the geometric sums $e^{-t\gamma/2}(1-e^{-t\gamma})^{-1}
=(2\sh(t\gamma/2))^{-1}$. As $t\to\infty$, $\cth\to1$ monotonically and the
thermal state converges to the vacuum in trace norm, giving the vacuum case.
\end{proof}

\begin{remark}[V0 acceptance]
\eqref{eq:latticefree}--\eqref{eq:Qt} are the formulas verified in V0: against
exact diagonalization to $6\times10^{-10}$ ($t=1$) and $2\times10^{-16}$
($t=2$), and against the independent position-space matrix-function evaluation
to $10^{-9}$ for $r_N\in\{8,64\}$, all $t$.
\end{remark}

\section{The jump representation and the exact mean shift (A2(a))}
\label{sec:jump}

This section proves the interacting marked-slice representation. Fix $t>0$ and
write $\omega(k):=\gamma_N(k)$ for the mode frequencies, $n:=2r_N$.

\subsection{The free kernel in Mehler form}

\begin{lemma}[Mehler kernel, verified]\label{lem:mehler}
In the rescaled coordinates $\bm X=\eps_N^{1/2}\bm\varphi$ and mode variables,
the free kernel factorizes over $k\in\Gamma_N$ into one-dimensional kernels
(for $H=\tfrac12(P^2+\omega^2X^2)$, which is $\gamma(a^*a+\tfrac12)$ per mode)
\begin{equation}
\begin{gathered}
 k^{\rm 1d}_\Delta(x,y)
 =\Bigl(\frac{\omega}{2\pi\sh(\omega\Delta)}\Bigr)^{1/2}
 \exp\Bigl[-\frac{A_\Delta}2(x^2+y^2)+B_\Delta\,xy\Bigr],\\
 A_\Delta=\omega\cth(\omega\Delta),\qquad
 B_\Delta=\frac{\omega}{\sh(\omega\Delta)} .
\end{gathered}
 \label{eq:mehler}
\end{equation}
\end{lemma}

\begin{proof}
Call the right side $\kappa_\Delta$. A direct computation gives
$\partial_\Delta\kappa_\Delta=\bigl(\tfrac12\partial_x^2
-\tfrac{\omega^2}2x^2\bigr)\kappa_\Delta$
(using $A'=-B^2$... explicitly: $A_\Delta'=-\omega^2/\sh^2$,
$B_\Delta'=-\omega^2\ch/\sh^2$, and the prefactor contributes
$-\tfrac12\omega\cth$; collecting terms on both sides matches coefficient by
coefficient in $\{1,x^2,y^2,xy\}$), and $\kappa_\Delta\to\delta$ as
$\Delta\downarrow0$ in the weak sense against $C_c$ (Gaussian concentration
with $A_\Delta-B_\Delta=\omega\tanh(\omega\Delta/2)\to0$,
$A_\Delta+B_\Delta\to\infty$). Hence for $\psi\in C_c^\infty$,
$u(\Delta):=\int\kappa_\Delta(\cdot,y)\psi(y)dy$ is a strongly differentiable
$L^2$-solution of $u'=-Hu$, $u(0)=\psi$, as is $e^{-\Delta H}\psi$; for the
difference $v$, $\tfrac d{d\Delta}\|v\|^2=-2\ip v{Hv}\le0$ and $v(0)=0$, so
$v\equiv0$. Thus $\kappa_\Delta$ is the kernel of $e^{-\Delta H}$.
\end{proof}

\subsection{The harmonic jump interpolant}

\begin{lemma}[jump interpolant]\label{lem:interp}
Fix $t>0$, a mode frequency $\omega>0$, and a jump $\rho\in\R$. There is
exactly one $H\in C^\infty([0,t])$ with
\[
 -\ddot H+\omega^2H=0,\qquad H(t)-H(0)=\rho,\qquad \dot H(t)=\dot H(0),
\]
namely $H(s)=A\ch(\omega s)+B\sh(\omega s)$ with
\begin{equation}
 A=-\frac\rho2,\qquad
 B=\frac\rho2\,\cth\Bigl(\frac{\omega t}2\Bigr).
 \label{eq:AB}
\end{equation}
It satisfies, for every $\Delta>0$ and $s\in[\Delta,t-\Delta]$, the
\emph{exact} three-term recursion
\begin{equation}
 H(s+\Delta)+H(s-\Delta)=2\ch(\omega\Delta)\,H(s),
 \label{eq:threeterm}
\end{equation}
its boundary values are $H(0)=-\rho/2$, $H(t)=+\rho/2$, and its classical
action is
\begin{equation}
 S(H):=\int_0^t\tfrac12\bigl(\dot H^2+\omega^2H^2\bigr)ds
 =\tfrac12\,\dot H(0)\,\rho
 =\frac{\omega\rho^2}4\,\cth\Bigl(\frac{\omega t}2\Bigr).
 \label{eq:action}
\end{equation}
\end{lemma}

\begin{proof}
The general solution of $-\ddot H+\omega^2H=0$ is
$A\ch(\omega s)+B\sh(\omega s)$. The two conditions read
$A(\ch(\omega t)-1)+B\sh(\omega t)=\rho$ and
$A\sh(\omega t)+B(\ch(\omega t)-1)=0$; solving,
$A=-B(\ch-1)/\sh$ and
$B[\sh^2-(\ch-1)^2]/\sh=\rho$, and since
$\sh^2-(\ch-1)^2=2(\ch-1)$,
$B=\rho\,\sh(\omega t)/(2(\ch(\omega t)-1))=\tfrac\rho2\cth(\omega t/2)$
and $A=-\rho/2$; uniqueness because the homogeneous system
($\rho=0$) forces $A=B=0$. \eqref{eq:threeterm} is the addition formula
$\ch(\omega(s\pm\Delta))$, $\sh(\omega(s\pm\Delta))$ summed. For
\eqref{eq:action}: integrate by parts,
$S(H)=\tfrac12[H\dot H]_0^t=\tfrac12(H(t)\dot H(t)-H(0)\dot H(0))
=\tfrac12\dot H(0)(H(t)-H(0))$ using $\dot H(t)=\dot H(0)$; and
$\dot H(0)=\omega B=\tfrac{\omega\rho}2\cth(\omega t/2)$.
\end{proof}

For the lattice, define the vector interpolant $\bm H_p(s)$ mode by mode: in
the rescaled coordinates the translation by $p$ is the shift by
$\eps_N^{1/2}p$, whose $k$-th Fourier component is
$\eps_N^{1/2}\widetilde p(k)$; let $H_k(s)$ be the interpolant of
Lemma~\ref{lem:interp} with $\omega=\gamma_N(k)$ and jump
$\rho_k=\eps_N^{1/2}\widetilde p(k)$ (complex $\rho$: apply the lemma to real
and imaginary parts; reality of the resulting position-space path follows from
$\rho_{-k}=\bar\rho_k$), and let $\bm H_p$ be the inverse transform. Then, by
\eqref{eq:action} and $\eps_N|\widetilde p(k)|^2=|\widehat p(k)|^2/(2r_N)$,
\begin{equation}
 S_N(\bm H_p):=\sum_{k\in\Gamma_N}S(H_k)
 =\frac1{4\cdot 2r_N}\sum_{k\in\Gamma_N}
 \gamma_N(k)\,\cth\Bigl(\frac{t\gamma_N(k)}2\Bigr)|\widehat p(k)|^2 ,
 \label{eq:SNp}
\end{equation}
which is exactly the momentum half of $\tfrac14Q^{(t)}_N$; and
$\bm H_p(0)=-p/2$, $\bm H_p(t)=+p/2$ in the original coordinates
(the mode relation $H_k(0)=-\rho_k/2$ is linear, so it inverts to
$\bm X$-shift $-\eps_N^{1/2}p/2$, i.e.\ $\bm\varphi$-shift $-p/2$).

\subsection{The exact change of variables}

Let $0=s_0<s_1<\cdots<s_M=t$ be any partition,
$\Delta_i:=s_{i+1}-s_i$. Define the (unnormalized) \emph{jump chain
integrals}: for bounded measurable $f_i:\R^n\to\C$,
\[
 J_p[f_0,\dots,f_{M-1}]
 :=\int\!\cdots\!\int
 \prod_{i=0}^{M-1}k^{(0)}_{\Delta_i}(\bm y_{i+1},\bm y_i)\,
 \prod_{i=0}^{M-1}f_i(\bm y_i)\;d\bm y_0\cdots d\bm y_{M-1}
 \Bigg|_{\bm y_M:=\bm y_0+p} ,
\]
and $J^{\rm per}$ the same with $\bm y_M:=\bm y_0$ (periodic chain).

\begin{lemma}[exact shift identity]\label{lem:shift}
For every partition and all bounded measurable $f_i$,
\begin{equation}
 J_p[f_0,\dots,f_{M-1}]
 =e^{-S_N(\bm H_p)}\;
 J^{\rm per}\bigl[f_0(\cdot+\bm H_p(s_0)),\dots,
 f_{M-1}(\cdot+\bm H_p(s_{M-1}))\bigr].
 \label{eq:shiftid}
\end{equation}
\end{lemma}

\begin{proof}
Work per mode in the rescaled Fourier coordinates (the free chain factorizes
by Lemma~\ref{lem:mehler}; the substitution below is a composition of
one-dimensional ones). Substitute $\bm y_i=\bm\eta_i+\bm H_p(s_i)$ (unit
Jacobian). Each Mehler factor transforms by
\[
 \log k^{\rm 1d}_{\Delta}(\eta'+a',\eta+a)
 =\log k^{\rm 1d}_{\Delta}(\eta',\eta)
 +\bigl[(-A_\Delta a'+B_\Delta a)\eta'
 +(-A_\Delta a+B_\Delta a')\eta\bigr]
 -S^{\rm cl}_\Delta(a',a),
\]
where $S^{\rm cl}_\Delta(a',a):=\tfrac{A_\Delta}2(a'^2+a^2)-B_\Delta a'a$
is the classical action of the harmonic path from $a$ to $a'$ in time
$\Delta$ (read off Lemma~\ref{lem:mehler}). Sum over the chain with
$a_i=H(s_i)$ (and $a_M=H(s_M)=H(s_0)+\rho$, matching
$\bm y_M=\bm y_0+p\mapsto\bm\eta_M=\bm\eta_0$, so the chain closes
periodically in $\bm\eta$). The linear terms at the node $\eta_i$
($1\le i\le M-1$) collect to
\[
 \bigl[-A_{\Delta_{i-1}}H(s_i)+B_{\Delta_{i-1}}H(s_{i-1})\bigr]
 +\bigl[-A_{\Delta_i}H(s_i)+B_{\Delta_i}H(s_{i+1})\bigr]=0 :
\]
indeed, for the \emph{harmonic} $H$, $H(s_{i\pm1})=\ch(\omega\Delta)H(s_i)
\pm\sh(\omega\Delta)\,\dot H(s_i)/\omega$ (addition formulas), so
$-A_\Delta H(s_i)+B_\Delta H(s_i\mp\Delta)
=\mp\,\dot H(s_i)/\!\ldots$ — explicitly,
$-\omega\cth(\omega\Delta)H(s_i)
+\tfrac{\omega}{\sh(\omega\Delta)}
[\ch(\omega\Delta)H(s_i)\mp\tfrac{\sh(\omega\Delta)}\omega\dot H(s_i)]
=\mp\dot H(s_i)$, and the two contributions ($-$ from the left factor,
$+$ from the right factor) cancel. At the glued node $\eta_0=\eta_M$ the same
computation applies with the one-sided derivatives $\dot H(0)$ and
$\dot H(t)$, which are equal by construction, so that node's linear term
vanishes as well. The constants collect to
$-\sum_iS^{\rm cl}_{\Delta_i}(H(s_{i+1}),H(s_i))$; since $H$ restricted to
$[s_i,s_{i+1}]$ \emph{is} the classical trajectory between its endpoints, each
term equals the true action integral
$\int_{s_i}^{s_{i+1}}\tfrac12(\dot H^2+\omega^2H^2)$, and the sum telescopes
to $S(H)$ of \eqref{eq:action}. Summing over modes gives
$e^{-S_N(\bm H_p)}$ and \eqref{eq:shiftid}.
\end{proof}

\begin{theorem}[A2(a): interacting mean-shift representation]\label{thm:meanshift}
For $t>0$, real $q,p$, and $\lambda\ge0$, with the periodic free ensemble
$\mu^{\rm per}_{N,t}$ (the Gaussian process on the time circle with the
covariance of the periodic Mehler chain; Remark~\ref{rem:ensemble}) and
$\mathcal V_N(\bm\varphi):=\lambda\eps_N\sum_x{:}\varphi_x^4{:}_{c_{N,L}}$
the \emph{interaction} multiplication function (the harmonic part of the
potential is carried by the ensemble; import (I1)), which satisfies the
pointwise bound
\[
 \mathcal V_N\ \ge\ -6\lambda\,\eps_N\,(2r_N)\,c_{N,L}^2
 \;=\;-12\lambda L\,c_{N,L}^2\;=\;-c_V
\]
of Proposition~\ref{prop:traceclass} (complete the square:
${:}u^4{:}_c=(u^2-3c)^2-6c^2$; this display is deliberately unnumbered, so
that the equation numbers cited by V2--V4 are unchanged):
\begin{equation}
 \frac{\Trace\bigl(e^{-tK_{N,L}}\,W_N(\eps_Nq+ip)\bigr)}
      {\Trace\bigl(e^{-tK_{N,L}}\bigr)}
 =e^{-S_N(\bm H_p)}\;
 \frac{\displaystyle\E_{\mu^{\rm per}_{N,t}}
 \Bigl[e^{\,i\eps_Nq\cdot\bm\eta(0)}\,
 e^{-\int_0^t\mathcal V_N(\bm\eta(s)+\bm H_p(s))\,ds}\Bigr]}
 {\displaystyle\E_{\mu^{\rm per}_{N,t}}
 \Bigl[e^{-\int_0^t\mathcal V_N(\bm\eta(s))\,ds}\Bigr]} ,
 \label{eq:meanshift}
\end{equation}
with $S_N(\bm H_p)$ as in \eqref{eq:SNp}; the CCR phase has cancelled against
the shift phase ($e^{i\theta}e^{i\eps_Nq\cdot\bm H_p(0)}
=e^{\frac i2\eps_Nq\cdot p}e^{-\frac i2\eps_Nq\cdot p}=1$). Expanding
$\mathcal V_N(\cdot+\bm H_p)$, the interaction changes by the finite Wick binomial
\begin{equation}
 \mathcal U_N(\bm\eta+\bm H_p)-\mathcal U_N(\bm\eta)
 =\lambda\eps_N\!\sum_x\!\int_0^t\!
 \Bigl[4H_x{:}\eta_x^3{:}+6H_x^2\,{:}\eta_x^2{:}
 +4H_x^3\,\eta_x+H_x^4\Bigr]ds,
 \label{eq:wickshift}
\end{equation}
(here $\mathcal U_N(\bm\eta):=\int_0^t\mathcal V_N(\bm\eta(s))\,ds$ denotes
the time-integrated interaction along a path), a polynomial with smooth
bounded deterministic coefficients
(by the closed form derived in the proof, for each fixed $N$,
$\|\bm H_p(s)\|_\infty\le\tfrac12\|p\|_{\ell^2(\Lambda_N)}$ uniformly in
$s\in[0,t]$ and $t\ge t_0$, with no dependence on $t_0$ or $m$; here
$\|p\|^2_{\ell^2(\Lambda_N)}=\sum_x|p(x)|^2$ is the counting-measure norm of the
model sheet (1.4), and the constant $\tfrac12$ is attained, since
$\bm H_p(0)=-p/2$).
This bound is \emph{not} uniform in $N$ along the $D4$ refinement: the mask is
orthogonal, $\sum_nh_nh_{n+2j}=\delta_{j0}$, so
$\|R_M^{M+r}p_M\|_{\ell^2(\Lambda_{M+r})}=2^{r/2}\|p_M\|_{\ell^2(\Lambda_M)}$ and
the right-hand side diverges. The $N$-uniform coarse-data bound that the route
actually consumes is $\sup_N\|\bm H_{p_N}\|_\infty\le\Sigma_M(p)$ of V2,
Lemma 6.3(i), which sums the $D4$ envelope over $\Gamma_N$ instead of applying
Cauchy--Schwarz over $2r_N$ modes; nothing downstream consumes the bound stated
here.
\end{theorem}

\begin{proof}
By Theorem~\ref{thm:kernel} and the Feynman--Kac formula \eqref{eq:FK}, the
numerator trace is
\[
 e^{i\theta}\!\int d\bm\varphi\,e^{i\eps_Nq\cdot\bm\varphi}\,
 k^{(0)}_t(\bm\varphi+p,\bm\varphi)\,
 \E^t_{\bm\varphi+p,\bm\varphi}\Bigl[e^{-\int_0^t\mathcal V_N}\Bigr]
 =e^{i\theta}\lim_{M\to\infty}
 J_p\Bigl[e^{i\eps_Nq\cdot(\cdot)}e^{-\Delta_0\mathcal V_N},
 e^{-\Delta_1\mathcal V_N},\dots\Bigr],
\]
where the second equality is the Riemann-sum approximation of the time
integral inside the bridge expectation: $\mathcal V_N$ is continuous, bounded below, and
the bridge paths are continuous, so
$\sum_i\Delta_i\mathcal V_N(B_{s_i})\to\int_0^t\mathcal V_N(B_s)ds$
pointwise on paths and the integrands are dominated by $e^{tc_V}$ times the
free chain (bounded-below $\mathcal V_N$, by the displayed pointwise bound
in Theorem~\ref{thm:meanshift}), allowing
dominated convergence under the (finite) free chain integral;
the chain form of the bridge-times-endpoint integral is the semigroup
(Chapman--Kolmogorov) property of $k^{(0)}$. Apply Lemma~\ref{lem:shift} at
each level $M$: the marked-slice functional $f_0$ produces the factor
$e^{i\eps_Nq\cdot(\bm\eta_0+\bm H_p(0))}$ with
$\bm H_p(0)=-p/2$, cancelling $e^{i\theta}$; the potential factors become
$e^{-\Delta_i\mathcal V_N(\bm\eta_i+\bm H_p(s_i))}$. Undo the Riemann sum on the periodic
side (same dominated convergence; $\bm H_p$ continuous), obtaining
$e^{-S_N(\bm H_p)}$ times the periodic FK integral with the shifted potential,
i.e.\ the numerator of \eqref{eq:meanshift}. The denominator is the case
$q=p=0$. Finally \eqref{eq:wickshift} is the Appell (binomial) property of Wick
powers: from ${:}u^4{:}_c=u^4-6cu^2+3c^2$,
\begin{gather*}
 {:}(\eta+H)^4{:}_c-{:}\eta^4{:}_c
 =\bigl[4H\eta^3+6H^2\eta^2+4H^3\eta+H^4\bigr]
 -6c\bigl[2H\eta+H^2\bigr]\\
 =4H{:}\eta^3{:}_c+6H^2{:}\eta^2{:}_c+4H^3\eta+H^4,
\end{gather*}
the last equality by $\eta^3={:}\eta^3{:}_c+3c\eta$,
$\eta^2={:}\eta^2{:}_c+c$, which cancels $-6c[2H\eta+H^2]$ identically (no
residual Wick-constant term survives in the Wick-power form); the uniform bound on
$\|\bm H_p\|_\infty$ follows from the \emph{closed form} of the interpolant,
not from \eqref{eq:AB} by the triangle inequality: the hyperbolic addition
formula turns \eqref{eq:AB} into
\[
 H_k(s)=\frac{\rho_k}2\,
 \frac{\sh\bigl(\gamma_N(k)(s-\tfrac t2)\bigr)}{\sh(\gamma_N(k)t/2)},
\]
a one-line verification: expanding
\[
 \sh\bigl(\gamma_N(s-\tfrac t2)\bigr)
 =\sh(\gamma_Ns)\ch\bigl(\tfrac{\gamma_Nt}2\bigr)
 -\ch(\gamma_Ns)\sh\bigl(\tfrac{\gamma_Nt}2\bigr)
\]
and dividing by $\sh(\gamma_Nt/2)$ returns exactly
$A_k\ch(\gamma_Ns)+B_k\sh(\gamma_Ns)$ with the $A_k,B_k$ of \eqref{eq:AB}. Hence $|H_k(s)|\le|\rho_k|/2$ for every $s\in[0,t]$, because
$|\sh|$ is even and increasing and $|s-\tfrac t2|\le\tfrac t2$. (The same
identity is recorded downstream as V2, Lemma 6.1; it is derived here from
\eqref{eq:AB} alone, so no import is created.) The triangle inequality applied to \eqref{eq:AB} would give only
$\tfrac{|\rho_k|}2\bigl(\ch(\gamma_Ns)
+\cth(\tfrac{t\gamma_N}2)\sh(\gamma_Ns)\bigr)$, which grows like
$e^{\gamma_Nt}$ and is useless in the regime $\gamma_N\sim2/\eps_N$; the
cancellation in the closed form is what makes the bound uniform, with no
$t_0$-dependence. The inverse transform is an $\ell^2\to\ell^\infty$ bounded
map on $2r_N$ modes with the stated $p$-dependence.
\end{proof}

\begin{remark}[the periodic ensemble]\label{rem:ensemble}
$\mu^{\rm per}_{N,t}$ is the finite product over $k\in\Gamma_N$ of the
one-dimensional Gaussian processes on the circle of circumference $t$ with
covariance
\begin{equation}
 c_\omega(s)=\frac{\ch\bigl(\omega(\tfrac t2-|s|_{\T_t})\bigr)}
 {2\omega\,\sh(\omega t/2)},\qquad \omega=\gamma_N(k),
 \label{eq:percov}
\end{equation}
which is the two-point function of the periodic Mehler chain: per mode, for
$0\le s\le t$ and $w=e^{-t\omega}$,
\[
 Z^{-1}\Trace\bigl(e^{-(t-s)H}Xe^{-sH}X\bigr)
 =\frac{e^{-s\omega}+we^{s\omega}}{2\omega\,(1-w)}=c_\omega(s),
\]
by the ladder computation
$\ip{n}{Xe^{-sH}X\,n}$-insertion and the geometric sums
$\sum(n{+}1)w^n=(1-w)^{-2}$, $\sum nw^n=w(1-w)^{-2}$. Existence of a
continuous version: $c_\omega(0)-c_\omega(s)\le C_\omega|s|$, so
$\E|\eta(s)-\eta(s')|^2\le C|s-s'|$ and, the process being Gaussian,
$\E|\eta(s)-\eta(s')|^4\le3C^2|s-s'|^2$; Kolmogorov--Chentsov applies. Its
finite-dimensional marginals are the normalized periodic Mehler chains (both
are Gaussian; the chain marginals are consistent by the semigroup property and
have the two-point function above, computed operator-side), which is the form
used in Lemma~\ref{lem:shift}.
To fix a normalization that the product description leaves implicit: the mode
variables carrying \eqref{eq:percov} are the unitary DFT modes of the
\emph{rescaled} coordinates $\bm X=\eps_N^{1/2}\bm\varphi$ of the
import note (I1) in \S\ref{sec:conv}, not of $\bm\varphi$ itself, whose DFT modes carry
$c_\omega/\eps_N$. Equivalently, in position space the $\bm\varphi$-process has
\[
 \E\bigl[\varphi_x(s)\,\varphi_y(s')\bigr]
 =\frac1{2L}\sum_{k\in\Gamma_N}e^{ik(x-y)}\,c_{\gamma_N(k)}(s-s'),
\]
which is the form matched verbatim by V2, Proposition 1.2(i); this was verified
against the explicitly diagonalized chain at $r_N=4,8,16$ to $4\times10^{-14}$.
\end{remark}

\begin{corollary}[free case closes the circle]\label{cor:freecircle}
For $\lambda=0$, \eqref{eq:meanshift} reads
\[
 F^{(0)}_N(t;\xi)=e^{-S_N(\bm H_p)}\,
 \E\bigl[e^{i\eps_Nq\cdot\bm\eta(0)}\bigr]
 =e^{-S_N(\bm H_p)}\,e^{-\frac{\eps_N^2}2\,q^\top C^{\rm per}_N(0)\,q},
\]
and by
\eqref{eq:percov} ($s=0$: $c_\omega(0)=\cth(\omega t/2)/(2\omega)$) and the
mode Parseval identity this equals
$\exp[-\tfrac14Q^{(t)}_N(\xi)]$: the momentum half is \eqref{eq:SNp} and the
position half is
\[
 \tfrac14\,(2r_N)^{-1}\textstyle\sum_k
 \cth(t\gamma_k/2)\,\gamma_k^{-1}|\widehat q(k)|^2
\]
--- in exact agreement with Theorem~\ref{thm:lattice-free}. The two proofs are
independent (operator ladder algebra vs.\ Gaussian chain calculus), so the
agreement is a genuine internal cross-check of \eqref{eq:meanshift}.
\end{corollary}

\section{The continuum free functional (A3(a), continuum)}\label{sec:cont}

Let $H_{0,L}=\mathrm d\Gamma(\mu)$, $\mu=(-\partial_x^2+m^2)^{1/2}$, on the
Fock space over $L^2(\T_L)$, with modes $k\in\Gamma_\infty=\tfrac\pi L\Z$ and
$\gamma(k)=(m^2+k^2)^{1/2}$. The model sheet's (2.13)--(2.16) fix
$\phi_{M,x}$, $R_M^\infty$, the scaling equation and $\beta_M$, but not the
continuum Weyl data themselves; we therefore fix them here, once, and use them
throughout Part~II. The one-particle space is
\[
 h_{\rm ct}:=\bigl\{\zeta=q+ip:\ \ip{q}{\gamma^{-1}q}+\ip{p}{\gamma p}<\infty\bigr\}
 =H^{-1/2}(\T_L)\oplus i\,H^{1/2}(\T_L),
\]
\emph{not} all of $L^2(\T_L;\C)$: for $p\in L^2\setminus H^{1/2}$ the operator
$\Pi(p)$ is not densely defined and $W_{\rm ct}(ip)$ is not unitary, so the larger
space would not carry a Weyl system. On $h_{\rm ct}$ the symplectic form
$\sigma_{\rm ct}(\zeta,\zeta')=\mathrm{Im}\int_{\T_L}\bar\zeta\zeta'$ is finite (it
pairs $H^{-1/2}$ with $H^{1/2}$), and we set
$W_{\rm ct}(q+ip):=e^{i(\Phi(q)+\Pi(p))}$ with the CCR phase
$\tfrac12\ip{q}{p}$, the continuum analogue of Proposition~\ref{prop:phase}. Note the
parametrisation is $\zeta=q+ip$, \emph{without} the lattice's $\eps_M$: with
$\eps_Mq$ in place of $q$ the Weyl relation would carry $\sigma_{\rm ct}/\eps_M$ and
$\beta_M$ would fail to be a $*$-homomorphism. The free vacuum
functional is
\[
 \omega^{(0)}_L\bigl(W_{\rm ct}(\zeta)\bigr)
 =\exp\Bigl[-\tfrac14\bigl(\ip{q}{\gamma^{-1}q}+\ip{p}{\gamma p}\bigr)\Bigr],
 \qquad\zeta=q+ip .
\]
With $q_\infty=\eps_M^{1/2}\sum_xq(x)\phi_{M,x}$ one has
$\int_{\T_L}q_\infty(y)e^{-iky}dy=\eps_M^{1/2}P_\infty(\eps_Mk)\widehat q_M(k)$
(using $\widehat\phi=P_\infty$, free-rate (1.2)), so this functional evaluated on
$R_M^\infty\xi$ is exactly the free-rate document's continuum vacuum
(1.7)--(1.8); the two sides of the gate are thus normalized identically, and
$R_M^\infty$ is symplectic by orthonormality of
$\{\phi_{M,x}\}_{x\in\Lambda_M}$. Finally $R_M^\infty$ maps into $h_{\rm ct}$:
$R_M^\infty\xi$ has Fourier coefficients $\eps_M^{1/2}P_\infty(\eps_Mk)\widehat q_M(k)$
with $|P_\infty(\ell)|\le C_D^{1/2}(1+|\ell|)^{-1-\eta/2}$ (free-rate Lemma 2.1), so
both $\ip{q_\infty}{\gamma^{-1}q_\infty}$ and $\ip{p_\infty}{\gamma p_\infty}$ converge
--- the latter needing $\sum_k\gamma(k)(1+\eps_M|k|)^{-2-\eta}<\infty$, which holds
since $2+\eta>2$. Hence every Weyl generator the gate names is defined, and
$\mathfrak W(\T_L)$ below is generated by $\{W_{\rm ct}(\zeta):\zeta\in h_{\rm ct}\}$.

\begin{theorem}[A3(a), continuum free]\label{thm:cont-free}
$e^{-tH_{0,L}}$ is trace class for every $t>0$, with
$\Trace e^{-tH_{0,L}}=\prod_{k\in\Gamma_\infty}(1-e^{-t\gamma(k)})^{-1}$
(convergent), and for every finite-mode symbol --- in particular for
$R_M^\infty\xi$, whose mode amplitudes $z_k$ are given by the continuum
analogue of Lemma~\ref{lem:amp} with $\sum_k|z_k|^2<\infty$ ---
\begin{equation}
 \frac{\Trace\bigl(e^{-tH_{0,L}}W_{\rm ct}(R_M^\infty\xi)\bigr)}
      {\Trace\,e^{-tH_{0,L}}}
 =\exp\Bigl[-\tfrac14\,Q^{(t)}_\infty(R_M^\infty\xi)\Bigr],
 \label{eq:contfree}
\end{equation}
where $Q^{(t)}_\infty$ is the $\cth(t\gamma(k)/2)$-weighted continuum
covariance form (the vacuum form of the free-rate document, eq.\ (1.8), with
the thermal weights inserted).
\end{theorem}

\begin{proof}
Here $H_{0,L}=\sum_k\gamma(k)\,a_k^*a_k$ (no zero-point subtlety: we use the
number-operator normalization; ratios are unaffected). The occupation basis
$\{|\bm n\rangle:\bm n\in\bigoplus_k\mathbb N_0\ \text{finitely supported}\}$
diagonalizes $e^{-tH_{0,L}}$ with eigenvalues $e^{-t\sum_k\gamma(k)n_k}$;
summability:
$\sum_{\bm n}e^{-t\sum\gamma n}
=\prod_k(1-e^{-t\gamma(k)})^{-1}$, finite because
$\sum_ke^{-t\gamma(k)}<\infty$ ($\gamma(k)\ge c|k|$). For a Weyl operator
with $\ell^2$ mode amplitudes, $W=\overline{\bigotimes}_kD_k(iz_k)$ in the
incomplete-tensor-product sense over the vacuum stabilization, and
$\ip{\bm n}{W\bm n}=\prod_k\ip{n_k}{D_k(iz_k)n_k}$: all but finitely many
factors are $\ip{0}{D(iz_k)0}=e^{-|z_k|^2/2}$, and the product converges
absolutely since $\sum|z_k|^2<\infty$. Then
\[
 \Trace(e^{-tH_{0,L}}W)
 =\sum_{\bm n}e^{-t\sum\gamma n}\prod_k\ip{n_k}{D_k(iz_k)n_k},
\]
absolutely convergent with majorant $\Trace e^{-tH_{0,L}}$
($|\ip{n}{Dn}|\le1$, unitarity), so Fubini--Tonelli for countable products
applies and the sum factorizes into
\[
 \prod_k\sum_{n_k}e^{-t\gamma n_k}\ip{n_k}{D_kn_k}.
\]
Lemma~\ref{lem:laguerre}
evaluates each factor; taking the ratio with the partition function and
inserting the continuum amplitude computation --- identical to
Lemma~\ref{lem:amp}, with the normalization
$(2r_N)^{-1}\sum_{\Gamma_N}$ replaced by
$\eps_M(2L)^{-1}\sum_{\Gamma_\infty}$ as in the free-rate document,
Lemma 1.2 --- gives \eqref{eq:contfree}.
\end{proof}

\section{The thermal fixed-coarse rate (A3(b))}\label{sec:rate}

Fix the coarse scale $M$, $\xi=\eps_Mq+ip$, and $t_0>0$. Write
$w_t(\gamma):=\cth(t\gamma/2)$, $\bar w:=\cth(t_0m/2)$, and
$x_0:=t_0m/2$. Define the pulled-back thermal forms $Q^{(t)}_{r,M}$ and
$Q^{(t)}_{\infty,M}$ exactly as in the free-rate document, Lemma 1.1/1.2, with
the weight $w_t(\gamma_{M+r}(k))$ resp.\ $w_t(\gamma(k))$ inserted (V0
verified the two implementations of the pullback agree to $2\times10^{-15}$,
i.e.\ at round-off level; \texttt{v0\_checkpoint.out}, item [2], records
$1.78\times10^{-15}$, and the value is not bit-reproducible across runs).

\begin{lemma}[derivative bounds for the thermal weights]\label{lem:wderiv}
Let $s_*:=\sup_{x\ge x_0}x/\sh^2(x)<\infty$. For all
$\gamma\ge m$ and $t\in[t_0,\infty]$:
\begin{align}
 \Bigl|\frac{d}{d\gamma}\bigl[\gamma\,w_t(\gamma)\bigr]\Bigr|
 &=\Bigl|\,w_t(\gamma)-\frac{t\gamma/2}{\sh^2(t\gamma/2)}\Bigr|
 \le\bar w+s_* ,
 \label{eq:Dp}\\
 \Bigl|\frac{d}{d\gamma}\bigl[\gamma^{-1}w_t(\gamma)\bigr]\Bigr|
 &=\Bigl|-\frac{w_t(\gamma)}{\gamma^2}
 -\frac{t/2}{\gamma\,\sh^2(t\gamma/2)}\Bigr|
 \le\frac{\bar w+s_*}{m^2} ,
 \label{eq:Dq}
\end{align}
and $1\le w_t(\gamma)\le\bar w$; at $t=\infty$, $w\equiv1$ and both
left-hand sides reduce to $1$ and $\gamma^{-2}\le m^{-2}$.
\end{lemma}

\begin{proof}
$\tfrac d{d\gamma}\cth(t\gamma/2)=-\tfrac t2\sh^{-2}(t\gamma/2)$, so the two
derivatives are as displayed. In \eqref{eq:Dp}, $w_t(\gamma)\le\bar w$
(monotonicity of $\cth$ in both arguments, $t\gamma\ge t_0m$) and
$\tfrac{t\gamma}2\sh^{-2}(\tfrac{t\gamma}2)=x/\sh^2x\big|_{x=t\gamma/2}\le
s_*$ since $x\ge x_0$. In \eqref{eq:Dq}, the first term is
$\le\bar w/m^2$ and the second is
$\tfrac{t/2}{\gamma\sh^2}=\tfrac{x}{\gamma^2\sh^2x}\le s_*/m^2$.
\end{proof}

\begin{theorem}[A3(b): thermal rate, uniform in $t$]\label{thm:thermalrate}
Let $\eta=2(1-\log_2\sqrt3)$, $\theta=2\eta/(5+\eta)$, and let
$r_*$, $J_r$, $C_D$, $A_q$, $A_p$, $b$, $\kappa$, $c_0$ be as in the free-rate
document, \S4. Then for all $r\ge r_*$ and all $t\in[t_0,\infty]$,
\begin{equation}
 \bigl|Q^{(t)}_{r,M}(\xi)-Q^{(t)}_{\infty,M}(\xi)\bigr|
 \le C^{\rm th}_Q\,2^{-\theta r},
 \qquad
 \Bigl|F^{(0)}_{M+r}\text{-side}-F^{(0)}_\infty\text{-side}\Bigr|
 \le\tfrac14\,C^{\rm th}_Q\,2^{-\theta r},
 \label{eq:thermalrate}
\end{equation}
where the second display compares the two thermal Weyl functionals of
Theorems \ref{thm:lattice-free} and \ref{thm:cont-free} on the embedded
generator, and
\begin{align}
 C^{\rm th}_Q&:=C^{\rm th}_{\rm low}\,(1+b^{-1})^{5}
 +\frac{2^{\eta}\,C^{\rm th}_*}{\pi\eta},
 \label{eq:CthQ}\\
 C^{\rm th}_{\rm low}
 &:=\frac{3\eps_M^3}{2L}
 \Bigl\{A_q^2\Bigl(\frac{(\bar w+s_*)\kappa^4}{24\,m^3}
 +\frac{c_0\bar w\,\kappa^2}{m}\Bigr)
 +A_p^2\Bigl(\frac{(\bar w+s_*)\kappa^4}{24\,m}
 +c_0\bar w\,(m+\kappa)\kappa^2\Bigr)\Bigr\},
 \nonumber\\
 C^{\rm th}_*&:=\bar w\,C_*
 =2\,\bar w\,C_D\Bigl(\frac{A_q^2}{m}+A_p^2\,(m+\eps_M^{-1})\Bigr).
 \nonumber
\end{align}
\end{theorem}

\begin{proof}
Follow the proof of the free-rate Theorem 4.1 step by step; only the weights
change. \emph{Low momenta} ($|n|\le J_r$): the two integrand errors are now,
by the triangle splits
\[
 |w(\gamma_N)\gamma_N|P_r|^2-w(\gamma)\gamma|P_\infty|^2|
 \le|w(\gamma_N)\gamma_N-w(\gamma)\gamma|\,|P_r|^2
 +w(\gamma)\gamma\,\bigl||P_r|^2-|P_\infty|^2\bigr|,
\]
and its $q$-analogue; the mean value theorem with Lemma~\ref{lem:wderiv} and
the imported dispersion bound $|\gamma_N-\gamma|\le\eps_N^2k^4/(24m)$
(free-rate Lemma 3.1) give
\[
 |w(\gamma_N)\gamma_N-w(\gamma)\gamma|
 \le(\bar w+s_*)\,\frac{\eps_N^2k^4}{24m},\qquad
 |w(\gamma_N)\gamma_N^{-1}-w(\gamma)\gamma^{-1}|
 \le\frac{\bar w+s_*}{m^2}\,\frac{\eps_N^2k^4}{24m},
\]
while the imported filter-tail bound (free-rate Lemma 2.2) is multiplied by
$w(\gamma)\gamma\le\bar w(m+\kappa(1+J))$ resp.\
$w(\gamma)/\gamma\le\bar w/m$. Summing over $|n|\le J_r$ exactly as in the
source (the counting bounds (4.16) are unchanged) yields
$E_{\rm low}\le4^{-r}C^{\rm th}_{\rm low}(1+J_r)^5$. \emph{High momenta}: the
envelope estimate (free-rate Lemma 2.1) is unchanged and the summand acquires
the factor $w\le\bar w$, so $C_*$ becomes $C^{\rm th}_*=\bar wC_*$ and the
integral-comparison tail (4.20)--(4.21) goes through verbatim. Balancing the
exponents as in the source ($2-5\alpha=\alpha\eta=\theta$) proves the first
display of \eqref{eq:thermalrate}; the second follows from
$|e^{-x/4}-e^{-y/4}|\le\tfrac14|x-y|$ ($x,y\ge0$). All bounds are uniform on
$t\in[t_0,\infty]$, with the $t=\infty$ case reducing to the printed vacuum
theorem (weights $\equiv1$, $s_*$-terms replaced by the original ones, and
$C^{\rm th}\ge C^{\rm vac}$).
\end{proof}

\begin{remark}[V0 acceptance]
The measured slopes at $t\in\{1,2,4\}$ and $t=\infty$ coincide to four
decimals (drift $0.0000$), consistent with the $t$-uniform bound
\eqref{eq:thermalrate}; the measured decay $2^{-1.011r}$ for the smooth V0
symbol is faster than the worst-case $2^{-\theta r}$, as expected.
\end{remark}

\section{A2(b): the continuum statement and its input clause}\label{sec:A2b}

\begin{definition}[A2(b), final form]\label{def:A2b}
Let $K_L$ be the continuum torus Hamiltonian in the A4b convention. The
continuum marked-slice representation is the statement: for $t>0$ and every
coarse symbol $\xi$,
\[
 \frac{\Trace\bigl(e^{-tK_L}W_{\rm ct}(R_M^\infty\xi)\bigr)}
      {\Trace\,e^{-tK_L}}
 =e^{-S_\infty(H_p)}\,
 \frac{\E_{\mu^{\rm per}_{t,L}}
 \bigl[e^{i\Phi(J_q)}e^{-\mathcal U(\eta+H_p)}\bigr]}
      {\E_{\mu^{\rm per}_{t,L}}\bigl[e^{-\mathcal U(\eta)}\bigr]},
\]
with $S_\infty$, $H_p$ the mode-wise objects of
Lemmas~\ref{lem:interp}--\ref{lem:shift} for the dispersion $\gamma(k)$ (the
free ingredients are proved: Theorem~\ref{thm:cont-free} and the
per-mode statements apply verbatim), $J_q$ the sharp-time position source
(admissible: Lemma 2.1 of the scoping note concerns momentum only), and
$\mathcal U$ the continuum interaction.
\end{definition}

\begin{remark}[input clause and amendment]\label{rem:amendment}
The only unproved ingredient in Definition~\ref{def:A2b} is the continuum
Feynman--Kac/trace formula for $e^{-tK_L}$ on the space--time torus (the
$\lambda=0$ case is Theorem~\ref{thm:cont-free}; the Gaussian shift lemma is
mode-wise and dimension-free). Accordingly, Statement A1 of the scoping note
is \emph{extended by the clause} (A1-iv): ``$e^{-tK_L}$ admits the periodic
Feynman--Kac representation over $\T_t\times\T_L$, compatible with cylinder
approximation.'' A2(b) is then a theorem given A1(i)--(iv), by the identical
chain of \S\ref{sec:jump}: kernel identity (the spectral/Fubini argument of
Theorem~\ref{thm:kernel} uses only trace-class, A1(ii)), shift identity
(mode-wise, proved), and Riemann/dominated convergence (uses
$e^{-\mathcal U}\in L^1$, part of A1(iv)'s package). The amendment is recorded
in \texttt{strategy.md}; A1 including (A1-iv) is the V3 deliverable, with
H{\o}egh-Krohn's finite-temperature construction as the import candidate for
precisely this clause.
\end{remark}

\section{Acceptance and claim ledger}\label{sec:ledger}

\begin{longtable}{P{5.6cm}P{2.6cm}P{5.6cm}}
\toprule
claim & status & ground \\
\midrule
\endhead
Schr\"odinger action of $W_N$, CCR phase $\tfrac12\eps_Nq\cdot p$ &
proved & Prop.~\ref{prop:phase}; V0 corner checks, residuals
$\le4\times10^{-15}$ \\
Trace class of $e^{-tK_{N,L}}$, $Z_N(t)\le e^{tc_V}Z^{(0)}_N(t)$ &
proved & Prop.~\ref{prop:traceclass} \\
Kernel identity for the marked-slice trace &
proved & Thm.~\ref{thm:kernel} (eigenfunction bounds, Fubini) \\
Free thermal functional, lattice (coth form) &
proved & Thm.~\ref{thm:lattice-free}; V0: ED residual $6\times10^{-10}$;
position-space cross-check passed at tolerance $10^{-9}$ \\
Single-mode thermal displacement &
proved & Lem.~\ref{lem:laguerre} (elementary Laguerre summation) \\
Harmonic jump interpolant; exact three-term recursion; action
$=\tfrac\omega4\cth(\tfrac{\omega t}2)\rho^2$ &
proved & Lem.~\ref{lem:interp} \\
Exact Gaussian shift identity at every partition &
proved & Lem.~\ref{lem:shift} (Mehler-chain harmonicity is exact) \\
Interacting mean-shift representation (lattice); phase cancellation;
Wick-binomial shift &
proved & Thm.~\ref{thm:meanshift} \\
Free case agreement of the two routes &
proved & Cor.~\ref{cor:freecircle} \\
Continuum free thermal functional &
proved & Thm.~\ref{thm:cont-free} (countable-product Fubini) \\
Thermal fixed-coarse rate, uniform on $[t_0,\infty]$ &
proved & Thm.~\ref{thm:thermalrate}; imports (I2); V0 slope drift $0.0000$ \\
A2(b) continuum interacting representation &
conditional & Def.~\ref{def:A2b} $+$ clause (A1-iv); deferred to V3
(Rem.~\ref{rem:amendment}) \\
\bottomrule
\end{longtable}

\noindent\emph{Phase V1 gate:} A2(a) and A3 are closed; A2(b) is in final
conditional form with its single input clause isolated and reassigned to V3.
The V0 acceptance criteria (formulas match the tables to rounding) are met by
Theorems \ref{thm:lattice-free}, \ref{thm:cont-free},
and \ref{thm:thermalrate}.

\end{document}
